% GENERATED FILE. Edit canonical lessons, then run npm run generate:books.
% canonical-source-hash: fnv1a64:6bddc21a67d0b391
% canonical-lessons: TA-C163-tayar-cey, TA-C163-parimaru, TA-C163-pakir, TA-C163-peram-pecu, TA-C163-caripar

\chapter{To prepare, To serve, To share, To bargain, To check}
\label{ch:ta-to-prepare-to-serve-to-share-to-bargain-to-check}
\hlchaptermodality{163}

\begin{chapteropening}
\textbf{By the end of this chapter:} \emph{I can say to prepare, to serve, to share, to bargain and to check.}

Complete the last lesson of chapter 163: I can say to prepare, to serve, to share, to bargain and to check.
\end{chapteropening}

\section[\texorpdfstring{\ta{தயார்} \ta{செய்} (tayār cey)}{தயார் செய் (tayār cey)}]{\ta{தயார்} \ta{செய்} (tayār cey) — to prepare}

\label{lesson:TA-C163-tayar-cey}



\begin{quote}
Before the new one: say the Tamil for to move, then the Tamil for to lift.
\end{quote}

\subsection*{You'll want to know: \ta{தயார்} \ta{செய்}}
\textbf{\ta{தயார்} \ta{செய்}} — \emph{tayār cey} — \textquotedblleft{}to prepare\textquotedblright{}.

\begin{grammarlens}[title={What you've built}]
One more word for this chapter.
\end{grammarlens}

\subsection*{Your turn}
\begin{itemize}
\raggedright
  \item \emph{Say it:} \emph{tayār cey}
  \item \emph{Say it:} \emph{tayār cey}, once more
  \item \emph{Say it:} say \emph{tayār cey}
  \item \emph{Recall:} say the Tamil for to move, then the Tamil for to lift, then say \emph{tayār cey} again
\end{itemize}

\begin{tcolorbox}[breakable,colback=teal!4,colframe=teal!35!black,title={Before you move on}]
What does \ta{தயார்} \ta{செய்} mean? (\textquotedblleft{}To prepare\textquotedblright{}.) Say it once more.
\end{tcolorbox}
\section[\texorpdfstring{\ta{பரிமாறு} (parimāṟu)}{பரிமாறு (parimāṟu)}]{\ta{பரிமாறு} (parimāṟu) — to serve (food)}

\label{lesson:TA-C163-parimaru}



\begin{quote}
Before the new one: say the Tamil for to lift, then the Tamil for to prepare.
\end{quote}

\subsection*{You'll want to know: \ta{பரிமாறு}}
\textbf{\ta{பரிமாறு}} — \emph{parimāṟu} — \textquotedblleft{}to serve (food)\textquotedblright{}.

\begin{grammarlens}[title={What you've built}]
One more word for this chapter.
\end{grammarlens}

\subsection*{Your turn}
\begin{itemize}
\raggedright
  \item \emph{Say it:} \emph{parimāṟu}
  \item \emph{Say it:} \emph{parimāṟu}, once more
  \item \emph{Say it:} say \emph{parimāṟu}
  \item \emph{Recall:} say the Tamil for to lift, then the Tamil for to prepare, then say \emph{parimāṟu} again
\end{itemize}

\begin{tcolorbox}[breakable,colback=teal!4,colframe=teal!35!black,title={Before you move on}]
What does \ta{பரிமாறு} mean? (\textquotedblleft{}To serve (food)\textquotedblright{}.) Say it once more.
\end{tcolorbox}
\section[\texorpdfstring{\ta{பகிர்} (pakir)}{பகிர் (pakir)}]{\ta{பகிர்} (pakir) — to share}

\label{lesson:TA-C163-pakir}



\begin{quote}
Before the new one: say the Tamil for to prepare, then the Tamil for to serve.
\end{quote}

\subsection*{You'll want to know: \ta{பகிர்}}
\textbf{\ta{பகிர்}} — \emph{pakir} — \textquotedblleft{}to share\textquotedblright{}.

\begin{grammarlens}[title={What you've built}]
One more word for this chapter.
\end{grammarlens}

\subsection*{Your turn}
\begin{itemize}
\raggedright
  \item \emph{Say it:} \emph{pakir}
  \item \emph{Say it:} \emph{pakir}, once more
  \item \emph{Say it:} say \emph{pakir}
  \item \emph{Recall:} say the Tamil for to prepare, then the Tamil for to serve, then say \emph{pakir} again
\end{itemize}

\begin{tcolorbox}[breakable,colback=teal!4,colframe=teal!35!black,title={Before you move on}]
What does \ta{பகிர்} mean? (\textquotedblleft{}To share\textquotedblright{}.) Say it once more.
\end{tcolorbox}
\section[\texorpdfstring{\ta{பேரம்} \ta{பேசு} (pēram pēcu)}{பேரம் பேசு (pēram pēcu)}]{\ta{பேரம்} \ta{பேசு} (pēram pēcu) — to bargain}

\label{lesson:TA-C163-peram-pecu}



\begin{quote}
Before the new one: say the Tamil for to serve, then the Tamil for to share.
\end{quote}

\subsection*{You'll want to know: \ta{பேரம்} \ta{பேசு}}
\textbf{\ta{பேரம்} \ta{பேசு}} — \emph{pēram pēcu} — \textquotedblleft{}to bargain\textquotedblright{}.

\begin{grammarlens}[title={What you've built}]
One more word for this chapter.
\end{grammarlens}

\subsection*{Your turn}
\begin{itemize}
\raggedright
  \item \emph{Say it:} \emph{pēram pēcu}
  \item \emph{Say it:} \emph{pēram pēcu}, once more
  \item \emph{Say it:} say \emph{pēram pēcu}
  \item \emph{Recall:} say the Tamil for to serve, then the Tamil for to share, then say \emph{pēram pēcu} again
\end{itemize}

\begin{tcolorbox}[breakable,colback=teal!4,colframe=teal!35!black,title={Before you move on}]
What does \ta{பேரம்} \ta{பேசு} mean? (\textquotedblleft{}To bargain\textquotedblright{}.) Say it once more.
\end{tcolorbox}
\section[\texorpdfstring{\ta{சரிபார்} (caripār)}{சரிபார் (caripār)}]{\ta{சரிபார்} (caripār) — to check}

\label{lesson:TA-C163-caripar}



\begin{quote}
Before the new one: say the Tamil for to share, then the Tamil for to bargain.
\end{quote}

\subsection*{You'll want to know: \ta{சரிபார்}}
\textbf{\ta{சரிபார்}} — \emph{caripār} — \textquotedblleft{}to check\textquotedblright{}.

\begin{grammarlens}[title={What you've built}]
One more word for this chapter.
\end{grammarlens}

\subsection*{Your turn}
\begin{itemize}
\raggedright
  \item \emph{Say it:} \emph{caripār}
  \item \emph{Say it:} \emph{caripār}, once more
  \item \emph{Say it:} say \emph{caripār}
  \item \emph{Recall:} say the Tamil for to share, then the Tamil for to bargain, then say \emph{caripār} again
\end{itemize}

\begin{tcolorbox}[breakable,colback=teal!4,colframe=teal!35!black,title={Before you move on}]
What does \ta{சரிபார்} mean? (\textquotedblleft{}To check\textquotedblright{}.) Say it once more.
\end{tcolorbox}
